\documentclass[11pt,a4paper]{article}
\usepackage[T1]{fontenc}
\usepackage[utf8]{inputenc}
\usepackage[finnish]{babel}
\usepackage{lmodern}
\usepackage[margin=2.3cm]{geometry}
\usepackage{amsmath,amssymb,amsthm,mathtools}
\usepackage{booktabs}
\usepackage{enumitem}
\usepackage{microtype}
\usepackage[hidelinks]{hyperref}

\newtheorem{lause}{Lause}
\newtheorem{lemma}{Lemma}
\newtheorem{propositio}{Propositio}
\newtheorem{seuraus}{Seuraus}
\theoremstyle{definition}
\newtheorem{maaritelma}{Määritelmä}
\newtheorem{oletus}{Oletus}
\theoremstyle{remark}
\newtheorem{huomautus}{Huomautus}

\newcommand{\E}{\mathbb{E}}
\renewcommand{\P}{\mathbb{P}}
\newcommand{\cC}{\mathcal{C}}
\newcommand{\cZ}{\mathcal{Z}}
\newcommand{\cT}{\mathcal{T}}
\newcommand{\cF}{\mathcal{F}}
\newcommand{\cB}{\mathcal{B}}
\newcommand{\cG}{\mathcal{G}}
\newcommand{\cH}{\mathcal{H}}
\newcommand{\cA}{\mathcal{A}}
\newcommand{\1}{\mathbf{1}}
\newcommand{\pol}[1]{\texttt{#1}}
\newcommand{\Bin}{\operatorname{Bin}}
\newcommand{\Poi}{\operatorname{Poisson}}
\newcommand{\Beta}{\operatorname{Beta}}
\DeclareMathOperator*{\argmax}{arg\,max}
% Taulukon rivit tiedostosta; primitiivinen \input toimii tabular-ympäristön sisällä.
\newcommand{\rivit}[1]{\csname @@input\endcsname #1 }

\title{Tärinähavaintoihin perustuvan hiekoituksen kohdistamisen hyöty\\[4pt]
\large Liukkaussimulaattorin matemaattinen malli ja todistus}
\author{Elias Wuolijoki\\ \small Junction Turku 2026, If-haaste}
\date{17.10.2026}

\begin{document}
\maketitle

\begin{abstract}
\noindent
Kuvaamme täsmällisesti tiedostossa \texttt{If/labra/v0/R/simuloi\_tammikuu.R}
toteutetun simulaattorin. Turun keskustan kulkuruuduilla on päivittäin vaihtuva
piilevä liukkaus, ohikulut ovat Poisson-jakautuneita, ja puhelimen tärinätunnistus
havaitsee liukastumisia virheineen. Huolto hiekoittaa rajallisen määrän alueita
kerrallaan, eikä päällekkäinen hiekoitus ole sallittu. Osoitamme, että
kun hiekoituksen vaikutus päättyy ennen seuraavaa huoltokierrosta,
tärinähavainnoista laskettuun beta-binomiaaliseen posterioriin perustuva
kohdistus tuottaa odotusarvoltaan enintään yhtä paljon kaatumisia kuin mikään
politiikka, joka ei käytä havaintoja. Parhaaseen niistä verrattuna ero on
aidosti positiivinen, kun tunnistus erottelee
liukastumiset virhehälytyksistä ja liukkaus vaihtelee ruuduittain. Kumpikin ehto
on myös välttämätön. Lisäksi näytämme, että paras havaintoja käyttämätön
politiikka hiekoittaa vilkkaimmat alueet. Tulokset koskevat simuloitua maailmaa;
Turun todellista liukkautta tai Ifin vahinkoja ei ole mitattu.
\end{abstract}

\tableofcontents

\section{Johdanto}

Palveluidean ytimessä on väite: jos kävelijöiden puhelimet tunnistavat
liukastumisia, hiekoitus voidaan ohjata paikkoihin, joissa siitä on eniten
hyötyä. Simulaattori on rakennettu tämän väitteen tutkimiseen. Tässä paperissa
väite muotoillaan lauseeksi ja todistetaan simulaattorin mallissa.

\paragraph{Tulokset lyhyesti.}
Olkoon $J(\pi)$ politiikan $\pi$ odotettu kaatumisten määrä tammikuussa.
Oletusten~\ref{ol:A} ja~\ref{ol:B} vallitessa
(hiekoituksen vaikutus päättyy ennen seuraavaa huoltokierrosta, ja vaikutusaika
on yhden vuorokauden sisällä) pätee
\[
  J(\pol{oraakkeli}) \;\le\; J(\pol{tuote}) \;<\; J(\pol{kulku})
  \;\le\; J(\pol{aikataulu}) \;\le\; J(\pol{ei\_huoltoa}).
\]
Keskimmäinen epäyhtälö on aito lauseen~\ref{la:aito} ehdoilla, jotka
simulaattorin oletusarvot täyttävät. Lause~\ref{la:sokea} osoittaa, että
\pol{kulku} on paras kaikista politiikoista, jotka eivät käytä
tärinähavaintoja. Siksi tulos ei koske vain aikataulun mukaista huoltoa.
Lause~\ref{la:valttamaton} näyttää, että ilman ruutujen välistä vaihtelua tai
ilman tunnistuksen erottelukykyä hyöty on täsmälleen nolla.

\paragraph{Mitä tässä ei väitetä.}
Lauseet koskevat mallia, jossa palvelu tuntee mallin rakenteen ja
parametrit. Ne eivät osoita, että Turun kävelyväylien liukkaus noudattaisi tätä
mallia, että hiekoituksen vaikutus olisi 80 prosenttia tai että tärinätunnistus
toimisi kuvatulla tarkkuudella. Hyödyn suuruus riippuu näistä oletuksista;
luvussa~\ref{sec:numeeriset} sitä arvioidaan simuloimalla.

\section{Malli}\label{sec:malli}

\subsection{Aika, paikka ja eksogeeniset suureet}

Aika etenee tunneittain, $t \in \cT = \{1, \dots, T\}$, missä $T = 744$
(1.--31.1.2026). Tunti $t = 1$ on 1.1.2026 klo 00--01 Suomen aikaa.
Funktio $d\colon \cT \to \{1, \dots, 31\}$ antaa tunnin paikallisen
vuorokauden. Se on kasvava, ja vuorokausi vaihtuu keskiyöllä.

Tarkastelualue on $300 \times 300$ metrin neliö VPK-talon ympärillä
(ETRS-TM35FIN). Se on jaettu $10 \times 10$ metrin ruutuihin. Malliin kuuluvat
ne $|\cC| = 85$ ruutua $c \in \cC$, joiden keskipiste on enintään 5 metrin
päässä jostakin neljästä käsin piirretystä reitistä. Huoltoalueet
$z \in \cZ$ ovat $50 \times 50$ metrin neliöitä, joissa on vähintään yksi
kulkuruutu; niitä on $|\cZ| = 20$. Kuvaus $\zeta\colon \cC \to \cZ$ kertoo
ruudun alueen, ja $\cC_z = \zeta^{-1}(z)$.

Seuraavat suureet ovat deterministisiä ja kaikkien politiikkojen tiedossa.
\begin{itemize}[leftmargin=*]
  \item \emph{Kulkemisen perustaso} $\lambda_c > 0$: niiden reittien
  tuntimäärien summa, joiden lähellä ruutu on. Alueen tasolla
  $\lambda(z) = \sum_{c \in \cC_z} \lambda_c$.
  \item \emph{Vuorokausirytmi}. Olkoon $h_t$ tunnin $t$ alkamiskellonaika Suomen
  aikaa. Tällöin
  \[
    r_t = \Bigl(0{,}1 + 0{,}8\, e^{-((h_t - 8)/2)^2} + 0{,}9\, e^{-((h_t - 17)/3)^2}\Bigr)
    \cdot \begin{cases} 0{,}7, & \text{lauantai tai sunnuntai},\\ 1, & \text{muulloin}.\end{cases}
  \]
  \item \emph{Sääkerroin}. Turku Artukaisen mitatuista tuntikeskiarvoista
  (lämpötila $T_t$, lumensyvyys $s_t$, sade $q_t$) lasketaan
  \[
    w_t = 0{,}7 + 1{,}2 \cdot \1\{-3 \le T_t \le 1\} + 0{,}6 \cdot \1\{s_t > 0\}
          + 0{,}5 \cdot \1\{q_t > 0,\ T_t < 1\} \in [0{,}7;\ 3{,}0].
  \]
  Liukastumisen perustodennäköisyys on $\pi_t = p_0 w_t$, missä $p_0 = 0{,}01$.
\end{itemize}
Sääkerroin on simulaattorin oma kokeellinen sääntö. Sitä ei ole kalibroitu
tienpinnan havaintoihin.

\subsection{Piilevä liukkaus}

\begin{maaritelma}[Liukkaustasot]\label{maar:tasot}
Olkoot $a, b > 0$, $G \in \mathbb{N}$, $B \sim \Beta(a,b)$ ja $F$ sen
kertymäfunktio. Asetetaan $q_j = F^{-1}(j/G)$, $j = 0, \dots, G$, ja
\[
  \ell_g = \E\bigl[B \mid q_{g-1} < B \le q_g\bigr], \qquad g = 1, \dots, G.
\]
Liukkauden jakauma on $\nu = \frac{1}{G} \sum_{g=1}^{G} \delta_{\ell_g}$, ja
sen odotusarvo on $\mu = \sum_g \nu_g \ell_g$.
\end{maaritelma}

\begin{lemma}\label{le:tasot}
Määritelmän~\ref{maar:tasot} tasoille pätee:
\begin{enumerate}[label=(\roman*)]
  \item $\ell_g = G \cdot \frac{a}{a+b} \bigl(I_{q_g}(a+1, b) - I_{q_{g-1}}(a+1, b)\bigr)$,
  missä $I_x$ on säännöllistetty epätäydellinen beetafunktio;
  \item $0 < \ell_1 < \ell_2 < \dots < \ell_G < 1$, kun $G \ge 1$;
  \item $\mu = a/(a+b)$.
\end{enumerate}
Erityisesti kun $b = 1$, on $F(x) = x^a$ ja
$\ell_g = \frac{G a}{a+1}\bigl((g/G)^{(a+1)/a} - ((g-1)/G)^{(a+1)/a}\bigr)$.
Oletusarvoilla $a = 1/4$, $b = 1$ saadaan
$\ell_g = \bigl(g^5 - (g-1)^5\bigr)/(5 G^4)$.
\end{lemma}

\begin{proof}
Olkoon $f_{a,b}$ beetajakauman tiheys. Koska $B(a+1, b)/B(a, b) = a/(a+b)$,
pätee $x f_{a,b}(x) = \frac{a}{a+b} f_{a+1,b}(x)$. Jokaisen välin
$(q_{g-1}, q_g]$ todennäköisyys on $1/G$, joten
\[
  \ell_g = G \int_{q_{g-1}}^{q_g} x f_{a,b}(x)\,dx
         = G \frac{a}{a+b} \int_{q_{g-1}}^{q_g} f_{a+1,b}(x)\,dx,
\]
mikä on (i). Tiheys on positiivinen välillä $(0,1)$, joten
$0 = q_0 < q_1 < \dots < q_G = 1$, ja ehdollinen odotusarvo on aidosti välin
sisällä: $q_{g-1} < \ell_g < q_g$. Tästä seuraa (ii). Kohta (iii) seuraa
kokonaisodotusarvon kaavasta:
$\mu = \sum_g \frac{1}{G} \E[B \mid \text{väli } g] = \E B$.
Kun $b = 1$, on $f(x) = a x^{a-1}$, $F(x) = x^a$ ja $q_j = (j/G)^{1/a}$, joten
$\ell_g = G \int_{q_{g-1}}^{q_g} a x^a\,dx = \frac{G a}{a+1}\bigl(q_g^{a+1} - q_{g-1}^{a+1}\bigr)$,
ja $q_g^{a+1} = (g/G)^{(a+1)/a}$. Arvolla $a = 1/4$ on $Ga/(a+1) = G/5$ ja $(a+1)/a = 5$.
\end{proof}

Oletusarvoilla ($a = 1/4$, $b = 1$, $G = 40$) on $\mu = 0{,}2$,
$\ell_1 \approx 7{,}8 \cdot 10^{-8}$ ja $\ell_{40} \approx 0{,}951$.
Jatkuvassa jakaumassa $\P(B < 0{,}004) = 0{,}004^{1/4} \approx 0{,}25$ ja
$\P(B > 0{,}5) = 1 - 0{,}5^{1/4} \approx 0{,}16$. Neljännes ruuduista on siis
käytännössä pitäviä, ja joka kuudes on hyvin liukas.

\begin{maaritelma}[Liukkauden dynamiikka]\label{maar:dynamiikka}
Olkoon $\rho \in [0, 1]$ pysyvyys. Jokaiselle ruudulle $c$ asetetaan
$L_{c,1} \sim \nu$. Kun $d \ge 1$, $L_{c,d+1} = L_{c,d}$ todennäköisyydellä
$\rho$, ja muuten $L_{c,d+1}$ arvotaan uudelleen jakaumasta $\nu$. Kaikki
arvonnat ovat riippumattomia ruutujen ja päivien välillä.
\end{maaritelma}

\begin{lemma}\label{le:markov}
Kaikilla $c$, $d$ ja $j \ge 0$ pätee $L_{c,d} \sim \nu$ ja
$\E[L_{c,d+j} \mid L_{c,1}, \dots, L_{c,d}] = \mu + \rho^j (L_{c,d} - \mu)$.
\end{lemma}

\begin{proof}
Siirtymäydin on $P(\ell_g, \ell_h) = \rho \1\{g = h\} + (1-\rho)\nu_h$, ja
$\nu P = \rho \nu + (1-\rho)\nu = \nu$. Koska alkujakauma on $\nu$, induktio
antaa $L_{c,d} \sim \nu$. Markovin ominaisuuden nojalla
$\E[L_{c,d+1} \mid L_{c,\le d}] = \rho L_{c,d} + (1-\rho)\mu$, ja väite seuraa
induktiolla luvun $j$ suhteen.
\end{proof}

Simulaattorin oletus on $\rho = 0$: ruudun liukkaus arvotaan joka päivä
uudelleen.

\subsection{Ohikulut, liukastumiset, kaatumiset ja havainnot}

\begin{maaritelma}[Todennäköisyysavaruus]\label{maar:avaruus}
Todennäköisyysavaruudella $(\Omega, \cA, \P)$ on seuraavat keskenään
riippumattomat satunnaissuureiden perheet:
\begin{enumerate}[label=(\alph*)]
  \item ohikulut $N_{c,t} \sim \Poi(\lambda_c r_t)$, riippumattomia kaikilla $(c,t)$;
  \item liukkauden arvonnat määritelmän~\ref{maar:dynamiikka} mukaisesti;
  \item ohikulkukohtaiset tasajakautuneet suureet $U_{c,t,i,j}$,
  $i \in \mathbb{N}$, $j \in \{1, 2, 3\}$;
  \item sokean politiikan mahdollinen satunnaistus $V$ (määritelmä~\ref{maar:politiikka}).
\end{enumerate}
Merkitään tunnin $t$ satunnaisuutta
$\Psi_t = (N_{c,t}, U_{c,t,i,j})_{c, i, j}$ ja koko liukkausprosessia $L$.
Tällöin $L, V, \Psi_1, \Psi_2, \dots$ ovat riippumattomia.
\end{maaritelma}

Olkoon $\gamma_{c,t} \in \{0,1\}$ tieto siitä, onko ruudun $c$ alue hiekoitettu
tunnilla $t$; se määräytyy politiikasta (luku~\ref{sec:huolto}). Olkoot
$e \in [0,1]$ hiekoituksen teho, $\varphi$ kaatumisen todennäköisyys
liukastumisesta, $\eta$ tunnistuksen herkkyys ja $\varepsilon$
virhehavainnon todennäköisyys. Liukastumisen todennäköisyys ohikulkua kohti on
\begin{equation}\label{eq:p}
  p_{c,t} = \pi_t\, L_{c,d(t)}\, (1 - e\,\gamma_{c,t}).
\end{equation}
Oletamme, että $\pi_t \ell_G < 1$ kaikilla $t$. Oletusarvoilla
$\max_t \pi_t \ell_G \le 0{,}03 \cdot 0{,}951$.

Ohikulku $i = 1, \dots, N_{c,t}$ johtaa liukastumiseen, jos
$U_{c,t,i,1} \le p_{c,t}$. Liukastuminen johtaa kaatumiseen, jos
$U_{c,t,i,2} \le \varphi$. Tärinätunnistus hälyttää, jos ohikulku on johtanut
liukastumiseen ja $U_{c,t,i,3} \le \eta$, tai jos liukastumista ei ole tapahtunut
ja $U_{c,t,i,3} \le \varepsilon$. Merkitään liukastumisten, kaatumisten ja
havaintojen määriä $S_{c,t}$, $F_{c,t}$ ja $D_{c,t}$.

\begin{lemma}[Havaintojen ja kaatumisten jakauma]\label{le:havainnot}
Olkoon $\cG_t = \sigma(L, V, \Psi_1, \dots, \Psi_{t-1}, N_{\cdot,t})$, ja olkoon
$\gamma_{\cdot,t}$ $\sigma(L, V, \Psi_{<t})$-mitallinen. Tällöin ehdolla $\cG_t$
suureet $D_{c,t}$, $c \in \cC$, ovat riippumattomia ja
\[
  D_{c,t} \sim \Bin\bigl(N_{c,t}, \theta_t(L_{c,d(t)}, \gamma_{c,t})\bigr),
  \qquad \theta_t(\ell, \gamma) = \varepsilon + (\eta - \varepsilon)\, \pi_t\, \ell\, (1 - e \gamma).
\]
Lisäksi $\E[F_{c,t} \mid \cG_t] = \varphi\, N_{c,t}\, p_{c,t}$.
\end{lemma}

\begin{proof}
Ehdolla $\cG_t$ suure $p_{c,t}$ on tunnettu, ja $U_{c,t,\cdot,\cdot}$ on
riippumaton $\cG_t$:stä. Ohikulun hälytyksen todennäköisyys on
$p_{c,t}\eta + (1 - p_{c,t})\varepsilon = \theta_t(L_{c,d(t)}, \gamma_{c,t})$,
ja ohikulut ovat riippumattomia. Samoin ohikulku johtaa kaatumiseen
todennäköisyydellä $p_{c,t}\varphi$.
\end{proof}

\begin{huomautus}[Toteutus ja yhteiset satunnaisluvut]\label{hu:crn}
Koodi arpoo summat $S$, $F$ ja $D$ suoraan binomijakaumista käänteisen
kertymäfunktion avulla. Sama tasajakautunut luku syötetään jokaiselle
politiikalle. Ehdolla $(N_{c,t}, p_{c,t})$ määrien yhteisjakauma on sama kuin
yllä. Jokaisen politiikan oma reunajakauma on siis täsmälleen tämä malli, ja
lauseet koskevat juuri sitä. Yhteiset satunnaisluvut kytkevät politiikat vain
pienentääkseen erotusten estimoinnin varianssia.
\end{huomautus}

\subsection{Huolto}\label{sec:huolto}

\begin{maaritelma}[Huoltokierrokset]\label{maar:huolto}
Olkoot $\Delta, H, K \in \mathbb{N}$: huoltoväli, vaikutusaika ja alueiden
määrä kierroksella. Kierrokset alkavat tunneilla
$\tau_k = 1 + (k-1)\Delta$, $k = 1, \dots, \kappa$, missä
$\kappa = \lceil T/\Delta \rceil$. Kierroksen vaikutusikkuna on
$W_k = \{\tau_k, \dots, \tau_k + H - 1\} \cap \cT$. Kierroksella valitaan
joukko $A_k \subseteq \cZ$, jossa $|A_k| = K$. Valinta saa tehdä vain kelpoisista
alueista:
\[
  A_k \subseteq E_k = \cZ \setminus \bigcup\bigl\{A_j : j < k,\ \tau_j + H > \tau_k\bigr\}.
\]
Tämä on \emph{päällekkäisen hiekoituksen kielto}. Hiekoitustieto on
$\gamma_{c,t} = \1\{\exists k : \zeta(c) \in A_k,\ t \in W_k\}$.
\end{maaritelma}

\begin{lemma}\label{le:erilliset}
Kiellon voimassa ollessa saman alueen vaikutusikkunat ovat erillisiä, joten
$\gamma_{c,t} = \sum_{k} \1\{\zeta(c) \in A_k\}\, \1\{t \in W_k\}$.
\end{lemma}

\begin{proof}
Olkoot $j < k$ ja $z \in A_j \cap A_k$. Kelpoisuuden mukaan
$\tau_j + H \le \tau_k$, joten $\max W_j < \tau_k = \min W_k$. Siksi summassa
on enintään yksi nollasta poikkeava termi.
\end{proof}

\subsection{Informaatio ja politiikat}

\begin{maaritelma}[Politiikka]\label{maar:politiikka}
Palvelun informaatio tunnin $t$ lopussa on
$\cF_t = \sigma(N_{c,s}, D_{c,s} : c \in \cC,\ s \le t)$, ja $\cF_0$ on triviaali.
\emph{Politiikka} on sääntö, jonka valinta $A_k$ on
$\cF_{\tau_k - 1}$-mitallinen ja kelpoinen. Koska $D$ riippuu aiemmista
valinnoista, $\cF_t = \cF^\pi_t$ riippuu politiikasta.
\begin{itemize}[leftmargin=*]
  \item Politiikka on \emph{sokea}, jos $A_k$ on mitallinen $\sigma$-algebran
  $\cB_{\tau_k - 1} \vee \sigma(V)$ suhteen, missä
  $\cB_t = \sigma(N_{c,s} : s \le t)$ ja $V$ on muusta riippumaton
  satunnaistus. Sokea politiikka ei siis käytä tärinähavaintoja.
  \item \emph{Oraakkelin} informaatio on
  $\cH_k = \sigma\bigl(L_{\cdot,d} : d \le d(\tau_k);\ \Psi_s : s < \tau_k\bigr)$.
  Se tuntee kuluvan päivän liukkauden mutta ei tulevaa satunnaisuutta.
  Jokaisella politiikalla $\cF^\pi_{\tau_k - 1} \subseteq \cH_k$.
\end{itemize}
\end{maaritelma}

\begin{maaritelma}[Tavoite]
Politiikan $\pi$ arvo on odotettu kaatumisten määrä
$J(\pi) = \E^\pi\bigl[\sum_{c \in \cC} \sum_{t \in \cT} F_{c,t}\bigr]$.
Simulaattorin sarake \texttt{expected\_falls} on
$\sum_{c,t} \varphi N_{c,t} p_{c,t}$. Lemman~\ref{le:havainnot} ja
torniominaisuuden nojalla sen odotusarvo on $J(\pi)$. Se on harhaton
estimaattori, jonka varianssi on pienempi kuin arvottujen kaatumisten.
\end{maaritelma}

\section{Beta-binomiaalinen posteriori}\label{sec:posteriori}

Olkoon politiikka deterministinen, eli $A_k$ on havaintohistorian
deterministinen funktio. Kaikki koodin politiikat ovat tällaisia, ja tasapelit
ratkaistaan alueen numeron mukaan.

\begin{propositio}[Suodin]\label{pr:suodin}
Ehdolla $\cF_t$ liukkausketjut $(L_{c,d})_d$, $c \in \cC$, ovat riippumattomia.
Jokaisen ruudun suodatusjakauma
$w_{c,g}(t) = \P(L_{c,d(t)} = \ell_g \mid \cF_t)$ saadaan rekursiolla.
\begin{enumerate}[label=(\roman*)]
  \item Alussa käytetään jakaumaa $\nu$. Päivän vaihtuessa jakauma ennustetaan
  seuraavalle päivälle: $w \mapsto \rho\, w + (1 - \rho)\nu$.
  \item Tunnin $t$ havainnot päivittävät jakaumaa:
  \[
    w_{c,g}(t) \propto w_{c,g}(t-1)\; \theta_t(\ell_g, \gamma_{c,t})^{D_{c,t}}
    \bigl(1 - \theta_t(\ell_g, \gamma_{c,t})\bigr)^{N_{c,t} - D_{c,t}}.
  \]
\end{enumerate}
Posteriorin odotusarvo on $m_{c}(t) = \sum_g w_{c,g}(t)\, \ell_g$.
\end{propositio}

\begin{proof}
Olkoon $h = (n_{c,s}, \delta_{c,s})_{c \in \cC,\, s \le t}$ havaintohistoria ja
$l = (l_{c,d})$ liukkauspolku. Politiikka on deterministinen, joten
$\gamma_{c,s} = \gamma_{c,s}(h_{<s})$. Lemman~\ref{le:havainnot} ja
peräkkäisen ehdollistamisen nojalla
\[
  \P(N_{\le t} = n, D_{\le t} = \delta \mid L = l)
  = \prod_{s \le t} \prod_{c} \Poi(n_{c,s}; \lambda_c r_s)
    \binom{n_{c,s}}{\delta_{c,s}} \theta_s\bigl(l_{c,d(s)}, \gamma_{c,s}(h_{<s})\bigr)^{\delta_{c,s}}
    \bigl(1 - \theta_s(\cdot)\bigr)^{n_{c,s} - \delta_{c,s}}.
\]
Polun $l$ funktiona tämä on tulo ruutukohtaisista tekijöistä, ja myös
priorijakauma on tulomuotoinen. Siksi posteriori on tulomuotoinen. Kunkin
ruudun posteriori on piilo-Markov-mallin posteriori. Siinä siirtymäydin on
lemman~\ref{le:markov} $P$ ja tunnin havainnon todennäköisyys on yllä oleva
binomitekijä. Eteenpäin-algoritmi antaa kohdat (i) ja (ii); ennuste
$wP = \rho w + (1 - \rho)\nu$.
\end{proof}

Koodissa tämä on funktio \texttt{paivita\_posteriori} ja päivän vaihdon rivi
\texttt{logw <- log(p\$pysyvyys * exp(logw) + (1 - p\$pysyvyys) / G)}.
Simulaattori arpoo $L$:n samasta diskreetistä jakaumasta $\nu$, joten laskettu
posteriori on tarkka, ei approksimaatio.

\begin{propositio}[Konjugaatti erikoistapaus]\label{pr:konjugaatti}
Olkoon priori jatkuva $\Beta(a,b)$, ja olkoon havaintojen todennäköisyys
$\theta = \ell$ eli $\eta = 1$, $\varepsilon = 0$ ja $\pi_s(1 - e\gamma) = 1$.
Jos saman päivän aikana on $n = \sum_s N_s$ ohikulkua ja $y = \sum_s D_s$
havaintoa, posteriori on $\Beta(a + y,\ b + n - y)$. Lisäksi
\[
  \E[L \mid \text{data}] = \frac{a + y}{a + b + n}
  = \frac{a + b}{a + b + n}\cdot\frac{a}{a+b} + \frac{n}{a + b + n}\cdot\frac{y}{n},
\]
ja $y$:n reunajakauma ehdolla $n$ on beta-binomijakauma $\mathrm{BB}(n, a, b)$.
\end{propositio}

\begin{proof}
Posteriorin tiheys on verrannollinen lausekkeeseen
$\ell^{a-1}(1-\ell)^{b-1} \cdot \ell^{y}(1-\ell)^{n-y}$, mikä on
$\Beta(a+y, b+n-y)$-tiheyden ydin. Reunajakauma on
$\binom{n}{y} B(a+y, b+n-y)/B(a,b)$, joka on beta-binomijakauman määritelmä.
\end{proof}

\paragraph{Intuitio.}
Priori $\Beta(a,b)$ vastaa $a + b$ kuviteltua ohikulkua, joista $a$ on
johtanut havaintoon. Posteriorin odotusarvo on painotettu keskiarvo:
vähällä datalla se pysyy lähellä priorin keskiarvoa ja paljolla datalla
lähestyy havaittua osuutta $y/n$. Simulaattorissa havainnon todennäköisyys on
liukkauden \emph{affiini} funktio,
$\theta = \varepsilon + (\eta - \varepsilon)\pi(1 - e\gamma)\,\ell$. Silloin
posteriori ei ole beetamuotoinen, ja se lasketaan siksi tarkasti tasojen
$\ell_1, \dots, \ell_G$ hilassa. Koodin tarkistus
\texttt{tarkista\_tammikuu.R} varmistaa, että hilalaskenta tuottaa
proposition~\ref{pr:konjugaatti} tuloksen, kun $G = 4000$.

\section{Hajotelma}\label{sec:hajotelma}

Määritellään alueen $z$ \emph{riskimassa} kierroksen $k$ vaikutusikkunassa:
\[
  R_{z,k} = \sum_{c \in \cC_z} \sum_{t \in W_k} N_{c,t}\, \pi_t\, L_{c,d(t)},
\]
sekä deterministinen \emph{altistus}
$\Lambda_{z,k} = \sum_{c \in \cC_z} \lambda_c \sum_{t \in W_k} r_t \pi_t
= \lambda(z)\, \kappa_k$, missä $\kappa_k = \sum_{t \in W_k} r_t \pi_t > 0$.
Riskimassa ei riipu politiikasta: se on ohikulkujen ja liukkauden funktio.

\begin{propositio}[Hyödyn hajotelma]\label{pr:hajotelma}
Olkoon $\cG_k$ mikä tahansa $\sigma$-algebra, jonka suhteen $A_k$ on mitallinen.
Jokaiselle päällekkäisen hiekoituksen kieltoa noudattavalle politiikalle pätee
\[
  J(\pi) = \mathsf{B} - e \varphi \sum_{k=1}^{\kappa}
  \E^\pi\Bigl[\sum_{z \in A_k} \E[R_{z,k} \mid \cG_k]\Bigr],
  \qquad \mathsf{B} = \varphi \mu \sum_{c,t} \lambda_c r_t \pi_t = J(\pol{ei\_huoltoa}).
\]
\end{propositio}

\begin{proof}
Lemman~\ref{le:havainnot} ja torniominaisuuden nojalla
$\E F_{c,t} = \varphi\,\E[N_{c,t}\, \pi_t L_{c,d(t)} (1 - e\gamma_{c,t})]$.
Lemman~\ref{le:erilliset} mukaan
\[
  \sum_{c,t} N_{c,t}\pi_t L_{c,d(t)} \gamma_{c,t}
  = \sum_k \sum_{z \in A_k} \sum_{c \in \cC_z} \sum_{t \in W_k} N_{c,t}\pi_t L_{c,d(t)}
  = \sum_k \sum_{z \in A_k} R_{z,k}.
\]
Koska $N \perp L$ ja $L_{c,d} \sim \nu$ (lemma~\ref{le:markov}), on
$\E[N_{c,t}\pi_t L_{c,d(t)}] = \lambda_c r_t \pi_t \mu$. Lopuksi
$\E[\1\{z \in A_k\} R_{z,k}] = \E[\1\{z \in A_k\}\E[R_{z,k} \mid \cG_k]]$,
koska $\1\{z \in A_k\}$ on $\cG_k$-mitallinen ja $R_{z,k} \ge 0$ on integroituva.
\end{proof}

\paragraph{Intuitio.}
Kielto takaa, että sama ruutu-tunti ei koskaan hyödy kahdesta hiekoituksesta.
Silloin politiikan hyöty on summa kierroksittaisista valinnoista: kuinka paljon
riskimassaa valitut alueet keskimäärin sisältävät valintahetken tiedon
valossa.

\begin{lemma}[Ehdollinen riskimassa]\label{le:ehdollinen}
Olkoon $\cG \subseteq \sigma(L, V, \Psi_{<\tau_k})$. Merkitään
$m_{c,d} = \E[L_{c,d} \mid \cG]$. Tällöin
$\E[R_{z,k} \mid \cG] = \sum_{c \in \cC_z} \lambda_c \sum_{t \in W_k} r_t \pi_t\, m_{c,d(t)}$.
Erityisesti:
\begin{enumerate}[label=(\roman*)]
  \item palvelulle ($\cG = \cF^\pi_{\tau_k - 1}$) $m_{c,d}$ on proposition~\ref{pr:suodin} posteriorin odotusarvo;
  \item sokealle politiikalle ($\cG = \cB_{\tau_k-1} \vee \sigma(V)$) $m_{c,d} = \mu$, joten $\E[R_{z,k} \mid \cG] = \mu \Lambda_{z,k}$;
  \item oraakkelille ($\cG = \cH_k$) $m_{c,d} = \mu + \rho^{d - d(\tau_k)}(L_{c,d(\tau_k)} - \mu)$, kun $d \ge d(\tau_k)$.
\end{enumerate}
\end{lemma}

\begin{proof}
Kun $t \ge \tau_k$, suure $N_{c,t}$ on osa satunnaisuutta $\Psi_t$, joten se
on riippumaton $\sigma$-algebrasta $\sigma(L, V, \Psi_{<\tau_k}) \supseteq \cG \vee \sigma(L)$.
Siksi $\E[N_{c,t} L_{c,d(t)} \mid \cG] = \E[N_{c,t}]\,\E[L_{c,d(t)} \mid \cG]$.
Kohdassa (ii) $L$ on riippumaton suureista $N$ ja $V$, joten
$\E[L_{c,d} \mid \cG] = \E L_{c,d} = \mu$. Kohta (iii) seuraa
lemmasta~\ref{le:markov}, koska tulevien päivien arvonnat ovat riippumattomia
$\sigma$-algebrasta $\cH_k$.
\end{proof}

\section{Päätulokset}\label{sec:lauseet}

\begin{oletus}[Ei siirtymää]\label{ol:A}
$H \le \Delta$.
\end{oletus}

\begin{oletus}[Ikkuna yhden vuorokauden sisällä]\label{ol:B}
Jokaisella $k$ ja jokaisella $t \in W_k$ pätee $d(t) = d(\tau_k)$.
\end{oletus}

Oletuksen~\ref{ol:A} vallitessa $\tau_j + H \le \tau_j + \Delta \le \tau_k$ aina,
kun $j < k$. Siksi $E_k = \cZ$, ja kaikki $K$ alueen joukot ovat kelpoisia.
Simulaattorin oletusarvot ovat $\Delta = H = 6$ ja $K = 2$. Koska 6 jakaa luvun
24, $\tau_1$ on keskiyöllä eikä tammikuussa ole kesäaikasiirtymää, myös
oletus~\ref{ol:B} on voimassa.

\begin{maaritelma}[Simulaattorin politiikat]\label{maar:koodi}
Merkitään $\cA_K = \{A \subseteq \cZ : |A| = K\}$.
\begin{itemize}[leftmargin=*]
  \item \pol{ei\_huoltoa}: ei hiekoitusta, $\gamma \equiv 0$.
  \item \pol{aikataulu}: $K$ seuraavaa kelpoista aluetta kiertävässä järjestyksessä.
  \item \pol{kulku}: $A_k \in \argmax_{A \in \cA_K} \sum_{z \in A} \Lambda_{z,k}$, eli $K$ aluetta, joilla $\lambda(z)$ on suurin.
  \item \pol{tuote}: $A_k \in \argmax_{A \in \cA_K} \sum_{z \in A} \hat S_{z,k}$, missä $\hat S_{z,k} = \E[R_{z,k} \mid \cF^{\pol{tuote}}_{\tau_k - 1}]$.
  \item \pol{oraakkeli}: $A_k \in \argmax_{A \in \cA_K} \sum_{z \in A} \E[R_{z,k} \mid \cH_k]$.
\end{itemize}
\end{maaritelma}

\begin{lause}[Paras sokea politiikka]\label{la:sokea}
Oletuksen~\ref{ol:A} vallitessa jokaiselle sokealle politiikalle $\beta$ pätee
\[
  J(\beta) \;\ge\; J^\ast_{\mathrm{sokea}} := \mathsf{B} - e\varphi\mu \sum_{k} \max_{A \in \cA_K} \sum_{z \in A} \Lambda_{z,k}
  \;=\; J(\pol{kulku}).
\]
\end{lause}

\begin{proof}
Valitaan propositiossa~\ref{pr:hajotelma} $\cG_k = \cB_{\tau_k - 1} \vee \sigma(V)$.
Lemman~\ref{le:ehdollinen} (ii) mukaan
$J(\beta) = \mathsf{B} - e\varphi\mu \sum_k \E\bigl[\sum_{z \in A_k} \Lambda_{z,k}\bigr]$.
Koska $A_k \in \cA_K$, summa on enintään maksimin suuruinen. Politiikka
\pol{kulku} saavuttaa maksimin jokaisella kierroksella. Koska
$\Lambda_{z,k} = \lambda(z)\kappa_k$ ja $\kappa_k > 0$, maksimoiva joukko
koostuu $K$ alueesta, joilla $\lambda(z)$ on suurin.
\end{proof}

\begin{seuraus}
$J(\pol{aikataulu}) \ge J(\pol{kulku})$. Sama pätee kaikkiin sääntöihin, jotka
valitsevat alueet kellonajan, sään tai havaittujen ohikulkujen perusteella.
\end{seuraus}

\paragraph{Intuitio.}
Ilman tärinähavaintoja jokaisen ruudun liukkauden paras arvio on keskiarvo
$\mu$. Hiekoituksen arvo on silloin verrannollinen pelkkään altistukseen.
Havaittujen ohikulkujen määrä ei kerro mitään liukkaudesta, koska ohikulut ja
liukkaus ovat mallissa riippumattomia.

\begin{lause}[Tärinähavaintojen arvo]\label{la:arvo}
Oletuksen~\ref{ol:A} vallitessa pätee $J(\pol{tuote}) \le J^\ast_{\mathrm{sokea}}$.
Siis $J(\pol{tuote}) \le J(\beta)$ jokaiselle sokealle politiikalle $\beta$.
Tarkemmin
\[
  J(\pol{kulku}) - J(\pol{tuote}) = e\varphi \sum_{k} \E\Bigl[\max_{A \in \cA_K} \sum_{z \in A} \hat S_{z,k}
  - \sum_{z \in A^{\pol{kulku}}_k} \hat S_{z,k}\Bigr] \;\ge\; 0,
\]
missä jokainen odotusarvon sisällä oleva termi on pisteittäin ei-negatiivinen.
\end{lause}

\begin{proof}
Valitaan propositiossa~\ref{pr:hajotelma} $\cG_k = \cF^{\pol{tuote}}_{\tau_k - 1}$:
\[
  J(\pol{tuote}) = \mathsf{B} - e\varphi \sum_k \E\Bigl[\max_{A \in \cA_K}\sum_{z \in A}\hat S_{z,k}\Bigr].
\]
Joukko $A^{\pol{kulku}}_k$ on deterministinen, ja
$\E \hat S_{z,k} = \E R_{z,k} = \mu \Lambda_{z,k}$. Siksi
$\E\bigl[\sum_{z \in A^{\pol{kulku}}_k} \hat S_{z,k}\bigr] = \mu \max_{A}\sum_{z \in A}\Lambda_{z,k}$.
Vähentämällä saadaan erotuksen kaava. Termi on ei-negatiivinen, koska maksimi
on vähintään sen arvo missä tahansa joukossa, myös joukossa $A^{\pol{kulku}}_k$.
\end{proof}

\paragraph{Intuitio.}
Tärinähavaintoja käyttävä palvelu voi aina valita saman joukon kuin paras
sokea politiikka. Se poikkeaa siitä vain, kun posteriori lupaa enemmän
riskimassaa. Posteriorin odotusarvo on keskimäärin priori, joten poikkeamat
eivät voi keskimäärin huonontaa valintaa. Kyseessä on epäyhtälö
$\E \max \ge \max \E$, informaation arvon perustulos.

\begin{lause}[Oraakkeli]\label{la:oraakkeli}
Oletuksen~\ref{ol:A} vallitessa $J(\pol{oraakkeli}) \le J(\pi)$ jokaiselle
politiikalle $\pi$, erityisesti $J(\pol{oraakkeli}) \le J(\pol{tuote})$.
\end{lause}

\begin{proof}
Koska $\cF^\pi_{\tau_k - 1} \subseteq \cH_k$, joukko $A^\pi_k$ on
$\cH_k$-mitallinen. Proposition~\ref{pr:hajotelma} nojalla
$J(\pi) = \mathsf{B} - e\varphi\sum_k \E\bigl[\sum_{z \in A^\pi_k}\E[R_{z,k}\mid\cH_k]\bigr]$.
Sisempi summa on enintään
$\max_{A \in \cA_K}\sum_{z \in A}\E[R_{z,k}\mid\cH_k]$, ja oraakkeli
saavuttaa tämän maksimin.
\end{proof}

\begin{lause}[Aito hyöty]\label{la:aito}
Oletetaan oletukset~\ref{ol:A} ja~\ref{ol:B} sekä seuraavat ehdot:
\begin{enumerate}[label=(C\arabic*)]
  \item $0 < \varepsilon < \eta \le 1$, $0 < e < 1$ ja $\varphi > 0$;
  \item $G \ge 2$;
  \item $K < |\cZ|$, $\Delta < 24$, ja on olemassa alueet
  $z^\ast \in A^{\pol{kulku}}_2$ ja $z' \notin A^{\pol{kulku}}_2$, joille
  $\ell_G\, \lambda(z') > \ell_1\, \lambda(z^\ast)$.
\end{enumerate}
Tällöin $J(\pol{tuote}) < J(\pol{kulku})$.
\end{lause}

\begin{proof}
Lauseen~\ref{la:arvo} kaavan termit ovat ei-negatiivisia. Riittää siis
osoittaa, että kierroksella $k = 2$ on positiivisella todennäköisyydellä
\begin{equation}\label{eq:aito}
  \max_{A \in \cA_K}\sum_{z \in A}\hat S_{z,2} > \sum_{z \in A^{\pol{kulku}}_2}\hat S_{z,2}.
\end{equation}
Kierros $\tau_2 = 1 + \Delta$ on ensimmäisenä päivänä, koska $\Delta < 24$.
Oletuksen~\ref{ol:B} ja lemman~\ref{le:ehdollinen} nojalla
$\hat S_{z,2} = \kappa_2 \sum_{c \in \cC_z}\lambda_c m_c$, missä
$m_c = \E[L_{c,1} \mid \cF_{\Delta}]$ ja $\kappa_2 > 0$. Jos
$\sum_{c \in \cC_{z'}}\lambda_c m_c > \sum_{c \in \cC_{z^\ast}}\lambda_c m_c$,
joukko $A^{\pol{kulku}}_2 \setminus\{z^\ast\} \cup \{z'\}$ antaa aidosti
suuremman summan, ja~\eqref{eq:aito} pätee.

Olkoon $n \in \mathbb{N}$ ja $E_n$ tapahtuma, jossa jokaisella ruudulla
$c \in \cC_{z'} \cup \cC_{z^\ast}$ on $N_{c,1} = n$ ja $N_{c,s} = 0$, kun
$2 \le s \le \Delta$. Lisäksi $D_{c,1} = n$, kun $c \in \cC_{z'}$, ja
$D_{c,1} = 0$, kun $c \in \cC_{z^\ast}$. Koska $\lambda_c r_s > 0$, ja
ehdon~(C1) sekä oletuksen $\pi_t\ell_G < 1$ nojalla
$0 < \varepsilon \le \theta_1(\ell, \gamma) < 1$, lemman~\ref{le:havainnot}
mukaan $\P(E_n) > 0$.

Tapahtumassa $E_n$ propositio~\ref{pr:suodin} antaa ensimmäisenä päivänä
(priori $\nu$, tunnit $s \ge 2$ eivät sisällä ohikulkuja)
\[
  w_{c,g} = \frac{\theta_1(\ell_g,\gamma_c)^n}{\sum_h \theta_1(\ell_h,\gamma_c)^n}
  \quad (c \in \cC_{z'}), \qquad
  w_{c,g} = \frac{(1-\theta_1(\ell_g,\gamma_c))^n}{\sum_h (1-\theta_1(\ell_h,\gamma_c))^n}
  \quad (c \in \cC_{z^\ast}).
\]
Ehtojen~(C1) ja~(C2) sekä lemman~\ref{le:tasot} (ii) nojalla $\theta_1(\ell, \gamma)$
on aidosti kasvava suureen $\ell$ suhteen kummallakin arvolla $\gamma \in \{0, 1\}$.
Siksi rajalla $n \to \infty$ painot keskittyvät tasolle $G$ alueella $z'$ ja
tasolle $1$ alueella $z^\ast$. Tällöin
$m_c \to \ell_G$, kun $c \in \cC_{z'}$, ja $m_c \to \ell_1$, kun $c \in \cC_{z^\ast}$.
Siis
\[
  \sum_{c \in \cC_{z'}}\lambda_c m_c - \sum_{c \in \cC_{z^\ast}}\lambda_c m_c
  \;\longrightarrow\; \ell_G\lambda(z') - \ell_1\lambda(z^\ast) > 0.
\]
Raja ei riipu muusta datasta, koska kunkin ruudun posteriori riippuu vain sen
omista havainnoista. Valitaan $n$ niin suureksi, että erotus on positiivinen
tapahtumassa $E_n$. Silloin~\eqref{eq:aito} pätee todennäköisyydellä vähintään
$\P(E_n) > 0$. Koska $e\varphi > 0$, väite seuraa.
\end{proof}

Simulaattorin oletusarvot täyttävät ehdot: $\varepsilon = 0{,}0002 < \eta = 0{,}8$,
$e = 0{,}8$, $\varphi = 0{,}02$, $G = 40$, $K = 2 < 20$ ja $\Delta = 6$.
Lisäksi $\ell_1/\ell_G \approx 8 \cdot 10^{-8}$, kun taas pienin $\lambda(z)$
on 16 ja suurin 157 ohikulkua tunnissa. Ehto (C3) on siis voimassa jokaiselle
parille.

\begin{lause}[Hyödyn välttämättömät ehdot]\label{la:valttamaton}
Oletuksen~\ref{ol:A} vallitessa:
\begin{enumerate}[label=(\roman*)]
  \item jos $\eta = \varepsilon$, niin $J(\pol{tuote}) = J(\pol{kulku})$;
  \item jos $G = 1$, niin $J(\pol{oraakkeli}) = J(\pol{tuote}) = J(\pol{kulku})$;
  \item jos $e = 0$ tai $\varphi = 0$, niin kaikilla politiikoilla $J(\pi) = \mathsf{B}$.
\end{enumerate}
\end{lause}

\begin{proof}
(i) Kun $\eta = \varepsilon$, $\theta_t$ ei riipu liukkaudesta, joten
proposition~\ref{pr:suodin} päivitys ei muuta jakaumaa. Silloin
$w(t) = \nu$ ja $m_c = \mu$, joten $\hat S_{z,k} = \mu\Lambda_{z,k}$ on
deterministinen. Lauseen~\ref{la:arvo} erotus on nolla.
(ii) Kun $G = 1$, $L \equiv \mu$, ja kaikkien kolmen politiikan pisteet ovat
$\mu\Lambda_{z,k}$.
(iii) Proposition~\ref{pr:hajotelma} hyötytermin kerroin on $e\varphi = 0$.
\end{proof}

\begin{seuraus}[Järjestys]
Oletuksen~\ref{ol:A} vallitessa
$J(\pol{oraakkeli}) \le J(\pol{tuote}) \le J(\pol{kulku}) \le J(\pol{aikataulu}) \le J(\pol{ei\_huoltoa}) = \mathsf{B}$.
Kun lauseen~\ref{la:aito} ehdot ovat voimassa, toinen epäyhtälö on aito.
\end{seuraus}

\begin{huomautus}[Pidempi vaikutusaika]\label{hu:H}
Kun $H > \Delta$, lemma~\ref{le:erilliset} ja propositio~\ref{pr:hajotelma}
pätevät edelleen kiellon ansiosta. Kelpoiset joukot $E_k$ kuitenkin riippuvat
aiemmista valinnoista. Lauseiden todistukset eivät siksi siirry sellaisinaan.
Havaintoihin perustuva valinta voi esimerkiksi viedä alueen pois käytöstä
hetkellä, jolloin sokea aikataulu olisi hyödyntänyt sitä. Koodi sallii tämän
asetelman, mutta sen tulokset ovat taulukossa~\ref{tab:herkkyys} vain empiirisiä.
\end{huomautus}

\begin{huomautus}[Ahneus]
\pol{tuote} maksimoi jokaisen kierroksen ehdollisen hyödyn erikseen. Lauseet
eivät väitä sitä optimaaliseksi kaikkien havaintoja käyttävien politiikkojen
joukossa. Hiekoitus pienentää alueen tulevien havaintojen
informatiivisuutta, koska $\theta$ pienenee kertoimella $(1-e)$. Tutkivampi
politiikka voisi siksi olla parempi. Lauseen~\ref{la:arvo} todistus kuitenkin
pätee sellaisenaan mille tahansa politiikalle, joka valitsee
$\argmax_A \sum_{z\in A}\E[R_{z,k} \mid \cG_k]$ jonkin $\sigma$-algebran
$\cG_k \subseteq \cH_k$ suhteen.
\end{huomautus}

\section{Koodin ja mallin vastaavuus}\label{sec:koodi}

\begin{table}[ht]
\centering
\small
\begin{tabular}{lll}
\toprule
Symboli & Koodissa (\texttt{tammikuu\_asetukset.R}, \texttt{R/simuloi\_tammikuu.R}) & Oletusarvo \\
\midrule
$a, b, G$ & \texttt{liukkaus\_a}, \texttt{liukkaus\_b}, \texttt{liukkaustasoja} & $0{,}25$; $1$; $40$ \\
$\rho$ & \texttt{pysyvyys} & $0$ \\
$p_0$ & \texttt{liukastuminen\_p} & $0{,}01$ \\
$\varphi$ & \texttt{kaatuminen\_liukastumisesta\_p} & $0{,}02$ \\
$\eta, \varepsilon$ & \texttt{tunnistus\_p}, \texttt{virhehavainto\_p} & $0{,}8$; $0{,}0002$ \\
$\Delta, H, K$ & \texttt{huoltovali\_h}, \texttt{huollon\_kesto\_h}, \texttt{huoltokohteita} & $6$; $6$; $2$ \\
$e$ & \texttt{huollon\_teho} & $0{,}8$ \\
$\ell_g$ & \texttt{liukkaustasot()} & lemma~\ref{le:tasot} \\
$w_{c,g}$ & \texttt{exp(logw)}, \texttt{paivita\_posteriori()} & propositio~\ref{pr:suodin} \\
$w_t$, $r_t$ & \texttt{saakerroin()}, \texttt{rytmi} & luku~\ref{sec:malli} \\
\bottomrule
\end{tabular}
\caption{Paperin merkinnät ja simulaattorin muuttujat.}
\label{tab:parametrit}
\end{table}

Koodi pisteyttää kierroksella $k$ alueen $z$ lausekkeella
\[
  \mathrm{piste}(z) = \sum_{c \in \cC_z}\lambda_c\bigl(\mu A_0 + (m_c - \mu)A_1\bigr),
  \qquad A_0 = \sum_{t \in W_k} r_t,\quad A_1 = \sum_{t \in W_k} r_t\,\rho^{d(t) - d(\tau_k)},
\]
missä $m_c = \mu$ (\pol{kulku}), posteriorin odotusarvo (\pol{tuote}) tai
$L_{c,d(\tau_k)}$ (\pol{oraakkeli}). Sääkerroin jätetään pois, jotta tulevaa
säätä ei tarvita.

\begin{propositio}
Oletusten~\ref{ol:A} ja~\ref{ol:B} vallitessa koodin \pol{kulku}, \pol{tuote}
ja \pol{oraakkeli} valitsevat määritelmän~\ref{maar:koodi} argmax-joukon.
\end{propositio}

\begin{proof}
Oletuksen~\ref{ol:B} vallitessa $A_1 = A_0$, joten
$\mathrm{piste}(z) = A_0\sum_{c \in \cC_z}\lambda_c m_c$. Lemman~\ref{le:ehdollinen}
mukaan vastaava tarkka arvo on $\kappa_k\sum_{c \in \cC_z}\lambda_c m_c$, missä
$m_c$ on sama ehdollinen odotusarvo. Proposition~\ref{pr:suodin} nojalla koodin
posteriori on tarkka. Koska $A_0 > 0$ ja $\kappa_k > 0$, pisteet ovat toistensa
positiivisia monikertoja, joten niiden argmax-joukot ovat samat. Oletuksen~\ref{ol:A}
vallitessa kelpoisuusrajoite ei poista yhtään joukkoa.
\end{proof}

Tarkistukset \texttt{tarkista\_tammikuu.R} ja \texttt{tarkista\_hyoty.R}
testaavat seuraavat asiat: diskretoinnin keskiarvon (lemma~\ref{le:tasot}),
suotimen vastaavuuden suoraan laskettuun posterioriin
(propositio~\ref{pr:suodin}), konjugaatin erikoistapauksen
(propositio~\ref{pr:konjugaatti}), päällekkäisen hiekoituksen kiellon myös
asetelmassa $H > \Delta$, tulevan sään vaikutuksen puuttumisen menneisiin
päätöksiin ja tapauksen $G = 1$ (lause~\ref{la:valttamaton}~(ii)).

\section{Numeeriset tulokset}\label{sec:numeeriset}

Lauseet kertovat tulosten järjestyksen mutta eivät erojen suuruutta. Erot
estimoidaan simuloimalla tammikuu 2026 useilla satunnaissiemenillä.
Jokaisella siemenellä on oma piilevä liukkaus ja omat ohikulut. Kaikki
politiikat käyttävät samoja satunnaislukuja (huomautus~\ref{hu:crn}).
Mittarina on \texttt{expected\_falls}. Luottamusvälit ovat muotoa
$\bar x \pm 1{,}96\, s/\sqrt{n}$ siementen yli; erotukset ovat parittaisia.
Komento: \texttt{Rscript paperi/laske\_taulukot.R}.

\begin{table}[ht]
\centering
\small
\begin{tabular}{lccc}
\toprule
Politiikka $\pi$ & $\hat J(\pi)$ [95 \% LV] & Vähenemä ei huoltoon, \% & $\hat J(\pol{kulku}) - \hat J(\pi)$ [95 \% LV] \\
\midrule
\rivit{taulukko_oletus.tex}
\bottomrule
\end{tabular}
\caption{Odotetut kaatumiset tammikuussa oletusasetuksilla (200 siementä).}
\label{tab:oletus}
\end{table}

\begin{table}[ht]
\centering
\footnotesize
\setlength{\tabcolsep}{3pt}
\begin{tabular}{lrrrrrrcrr}
\toprule
Asetelma & $n$ & \pol{ei} & \pol{aikat.} & \pol{kulku} & \pol{tuote} & \pol{oraak.} & \pol{kulku}$-$\pol{tuote} [95 \% LV] & $>0$ & Osuus \\
\midrule
\rivit{taulukko_herkkyys.tex}
\bottomrule
\end{tabular}
\caption{Herkkyys. Sarake ${>}0$ kertoo, monellako siemenellä \pol{tuote} oli
\pol{kulku}a parempi. Osuus on $(\hat J(\pol{kulku}) - \hat J(\pol{tuote}))/(\hat J(\pol{kulku}) - \hat J(\pol{oraakkeli}))$:
kuinka suuren osan täydellisen tiedon edusta tärinähavainnot saavuttavat.}
\label{tab:herkkyys}
\end{table}

% Tulkinta taulukoille, jotka tuotettiin komennolla Rscript paperi/laske_taulukot.R.
\paragraph{Tulkinta.}
Taulukoiden järjestys vastaa lauseita. Oletusasetuksilla
\pol{tuote} oli \pol{kulku}a parempi 193 siemenellä 200:sta, ja parittaisen
erotuksen luottamusväli $[0{,}38;\ 0{,}44]$ ei sisällä nollaa
(lause~\ref{la:aito}). Asetelmissa $G = 1$ ja $\eta = \varepsilon$ erotus oli
jokaisella siemenellä täsmälleen nolla, kuten lause~\ref{la:valttamaton}
ennustaa.

Suurin osa \pol{tuote}en edusta \pol{aikataulu}un nähden tulee kulkemisen
huomioimisesta: \pol{kulku} vähentää odotettuja kaatumisia noin 2,1 ja
tärinähavainnot vielä noin 0,4. Tärinähavainnot saavuttavat noin neljänneksen
täydellisen tiedon edusta. Oletusarvolla $\rho = 0$ liukkaus arvotaan
uudelleen keskiyöllä. Klo 00 kierroksella kuluvasta päivästä ei siksi ole
lainkaan tietoa, ja klo 06 kierroksella vain hiljaisen yön verran.
Informatiivisia ovat lähinnä klo 12 ja 18 kierrokset. Kun liukkaus säilyy
päivästä toiseen ($\rho = 0{,}9$), tieto kertyy, ja osuus nousee noin
60 prosenttiin.

Myös oraakkeli vähentää kaatumisia vain noin 24 prosenttia, koska huolto
kattaa kerrallaan 2 aluetta 20:stä. Asetelma $H = 12$ on oletuksen~\ref{ol:A}
ulkopuolella. Siinä \pol{tuote} oli edelleen keskimäärin parempi, mutta
lauseet eivät takaa tätä (huomautus~\ref{hu:H}).

Luvut ovat simuloituja kaatumisia $300 \times 300$ metrin alueella yhden kuukauden
aikana. Ne eivät ole vakuutusvahinkoja tai Turun todellisia tapahtumia.


\section{Rajoitukset}

\begin{itemize}[leftmargin=*]
  \item \emph{Malli on tunnettu.} Palvelun posteriori käyttää simulaattorin
  todellisia parametreja. Väärin määritelty priori tai väärä tunnistustarkkuus
  voi heikentää tulosta. Lause~\ref{la:arvo} ei kata tätä tapausta.
  \item \emph{Kaikki ohikulut havaitaan.} Sovelluksen käyttäjäosuus on
  poistettu. Osittainen otos pienentäisi $n$:ää ja siten informaatiota.
  Lauseet pätevät edelleen, jos havaitsematon osa on riippumaton liukkaudesta.
  \item \emph{Riippumattomuus.} Ohikulut, liukastumiset ja hälytykset ovat
  ehdollisesti riippumattomia. Ihmisten väistely, ryhmät ja saman henkilön
  toistuvat kulut puuttuvat.
  \item \emph{Liukkaus on vuorokauden ajan vakio}, ja ruudut ovat keskenään
  riippumattomia. Spatiaalinen korrelaatio parantaisi todennäköisesti
  ennustetta, koska naapurit lainaisivat tietoa toisilleen.
  \item \emph{Sääkerroin, hiekoituksen teho ja kaatumisosuus ovat oletuksia.}
  Niitä ei ole kalibroitu. Kaatumismäärät eivät ole vakuutusvahinkoja;
  muunnos euroiksi vaatii erillisen mallin.
  \item \emph{Reitit ja kulkemismäärät on piirretty käsin}, eivätkä ne ole
  mitattua aineistoa.
\end{itemize}

\end{document}
